% TODO make hw stylesheet
\documentclass[11pt]{scrartcl}
\usepackage[a4paper, total={6in, 8in}]{geometry}
\usepackage[usenames,dvipsnames,svgnames]{xcolor}
\usepackage[shortlabels]{enumitem}
\usepackage[framemethod=TikZ]{mdframed}
\usepackage{amsmath,amssymb,amsthm}
\usepackage[colorlinks]{hyperref}
\usepackage{multicol}
\usepackage{tabularx}

\hypersetup{
	colorlinks=true,
	linkcolor=blue,
	filecolor=magenta,      
	urlcolor=magenta,
	pdftitle={MAT 320 HW}
}

\usepackage{microtype}
\usepackage{mathtools}
\usepackage[headsepline]{scrlayer-scrpage}
\usepackage{thmtools}
\usepackage[textsize=footnotesize]{todonotes}

\mdfdefinestyle{mdbluebox}{roundcorner=10pt,innerbottommargin=9pt,
	linecolor=blue,backgroundcolor=TealBlue!5,}
\declaretheoremstyle[headfont=\sffamily\bfseries\color{MidnightBlue},
mdframed={style=mdbluebox},]{thmbluebox}

\mdfdefinestyle{mdbloobox}{frametitlefont=\bfseries,innerbottommargin=8pt,
	nobreak=true,backgroundcolor=TealBlue!5,linecolor=blue,}
\declaretheoremstyle[headfont=\bfseries\color{MidnightBlue},
mdframed={style=mdbloobox},headpunct={\\[3pt]},postheadspace=0pt,]{thmbloobox}

\mdfdefinestyle{mdredbox}{frametitlefont=\bfseries,innerbottommargin=8pt,
	nobreak=true,backgroundcolor=Salmon!5,linecolor=RawSienna,}
\declaretheoremstyle[headfont=\bfseries\color{RawSienna},
mdframed={style=mdredbox},headpunct={\\[3pt]},postheadspace=0pt,]{thmredbox}

\mdfdefinestyle{mdgreenbox}{linecolor=ForestGreen,backgroundcolor=ForestGreen!5,
	linewidth=2pt,rightline=false,leftline=true,topline=false,bottomline=false,}
\declaretheoremstyle[headfont=\bfseries\sffamily\color{ForestGreen!70!black},
mdframed={style=mdgreenbox},headpunct={ --- },]{thmgreenbox}

\mdfdefinestyle{mdorangebox}{linecolor=RedOrange,backgroundcolor=RedOrange!5,
	linewidth=2pt,rightline=false,leftline=true,topline=false,bottomline=false,}
\declaretheoremstyle[headfont=\bfseries\sffamily\color{RedOrange!70!black},
mdframed={style=mdorangebox},headpunct={ --- },]{thmorangebox}

\mdfdefinestyle{mdblackbox}{linecolor=black,backgroundcolor=RedViolet!5!gray!5,
	linewidth=3pt,rightline=false,leftline=true,topline=false,bottomline=false,}
\declaretheoremstyle[mdframed={style=mdblackbox}]{thmblackbox}

\declaretheoremstyle[spaceabove=6pt,spacebelow=6pt]{basehead}

\declaretheorem[style=basehead,name=Problem]{problem}
\declaretheorem[style=thmredbox,name=Problem,numbered=no]{problem*}
\declaretheorem[style=thmbloobox,name=Setup]{setup} % for config geo
\declaretheorem[style=thmredbox,name=Example,sibling=problem]{example}
\declaretheorem[style=thmredbox,name=$\bigstar$ Problem,sibling=problem]{niceprob}
\declaretheorem[style=thmbluebox,name=Theorem,sibling=problem]{theorem}
\declaretheorem[style=thmbluebox,name=Corollary,sibling=theorem]{corollary}
\declaretheorem[style=thmbluebox,name=Proposition,sibling=theorem]{prop}
\declaretheorem[style=thmbluebox,name=Lemma,numberwithin=problem]{lemma}
\declaretheorem[style=thmbluebox,name=Theorem,numbered=no]{theorem*}
\declaretheorem[style=thmbluebox,name=Lemma,numbered=no]{lemma*}
\declaretheorem[style=thmgreenbox,name=Claim,numberwithin=problem]{claim}
\declaretheorem[style=thmgreenbox,name=Claim,numbered=no]{claim*}
\declaretheorem[style=thmblackbox,name=Remark,numberwithin=problem]{remark}
\declaretheorem[style=thmblackbox,name=Remark,numbered=no]{remark*}
\declaretheorem[style=thmgreenbox,name=Definition,sibling=theorem]{definition}
\declaretheorem[style=thmgreenbox,name=Definition,numbered=no]{definition*}
\declaretheorem[style=thmorangebox,name=Warning,sibling=theorem]{warning}
\declaretheorem[style=thmorangebox,name=Warning,numbered=no]{warning*}
\declaretheorem[style=thmorangebox,name=Note,numbered=no]{note}
\declaretheorem[style=thmblackbox,name=Exercise,sibling=problem]{exercise}
\declaretheorem[style=thmblackbox,name=Exercise,numbered=no]{exercise*}
\declaretheorem[style=thmblackbox,name=Question,sibling=problem]{question}
\declaretheorem[style=thmblackbox,name=Question,numbered=no]{question*}

\addtolength{\textheight}{3.14cm}
\providecommand{\ol}{\overline}
\providecommand{\ul}{\underline}
\providecommand{\eps}{\varepsilon}
\providecommand{\half}{\frac{1}{2}}
\providecommand{\dang}{\measuredangle} %% Directed angle
\providecommand{\CC}{\mathbb C}
\providecommand{\FF}{\mathbb F}
\providecommand{\NN}{\mathbb N}
\providecommand{\QQ}{\mathbb Q}
\providecommand{\RR}{\mathbb R}
\providecommand{\ZZ}{\mathbb Z}
\providecommand{\dg}{^\circ}
\providecommand{\ii}{\item}
\providecommand{\setdiff}{\smallsetminus}
\providecommand{\alert}[1]{{\sffamily\textbf{\textcolor{blue}{#1}}}}
\providecommand{\inv}{^{-1}}
\providecommand{\mb}[1]{\mathbf{#1}}

\reversemarginpar

\newenvironment{soln}{\begin{proof}[Solution]}{\end{proof}}
\newenvironment{walkthrough}{\noindent \textbf{Walkthrough.}}{\medskip}
\newlist{walk}{enumerate}{3}
\setlist[walk]{label=\bfseries (\alph*)}

\providecommand{\pow}{\operatorname{Pow}}
\providecommand{\vocab}[1]{{\textbf{\textcolor{ForestGreen}{#1}}}}
\providecommand{\lcm}{\operatorname{lcm}}
\providecommand{\abs}[1]{\left|#1\right|}

\usepackage{asymptote}
\begin{asydef}
	defaultpen(fontsize(10pt));
	unitsize(1cm);
	size(8cm); // set a reasonable default
	usepackage("amsmath");
	usepackage("amssymb");
	settings.tex="pdflatex";
	settings.outformat="pdf";
	// Replacement for olympiad+cse5 which is not standard
	import geometry;
	// recalibrate fill and filldraw for conics
	void filldraw(picture pic = currentpicture, conic g, pen fillpen=defaultpen, pen drawpen=defaultpen)
	{ filldraw(pic, (path) g, fillpen, drawpen); }
	void fill(picture pic = currentpicture, conic g, pen p=defaultpen)
	{ filldraw(pic, (path) g, p); }
	// some geometry
	pair foot(pair P, pair A, pair B) { return foot(triangle(A,B,P).VC); }
	pair orthocenter(pair A, pair B, pair C) { return orthocentercenter(A,B,C); }
	pair centroid(pair A, pair B, pair C) { return (A+B+C)/3; }
	// cse5 abbreviations
	path CP(pair P, pair A) { return circle(P, abs(A-P)); }
	path CR(pair P, real r) { return circle(P, r); }
	pair IP(path p, path q) { return intersectionpoints(p,q)[0]; }
	pair OP(path p, path q) { return intersectionpoints(p,q)[1]; }
	path Line(pair A, pair B, real a=0.6, real b=a) { return (a*(A-B)+A)--(b*(B-A)+B); }
	// cse5 more useful functions
	picture CC() {
		picture p=rotate(0)*currentpicture;
		currentpicture.erase();
		return p;
	}
	pair MP(Label s, pair A, pair B = plain.S, pen p = defaultpen) {
		Label L = s;
		L.s = "$"+s.s+"$";
		label(L, A, B, p);
		return A;
	}
	pair Drawing(Label s = "", pair A, pair B = plain.S, pen p = defaultpen) {
		dot(MP(s, A, B, p), p);
		return A;
	}
	path Drawing(path g, pen p = defaultpen, arrowbar ar = None) {
		draw(g, p, ar);
		return g;
	}
\end{asydef}

\usepackage{mathrsfs}

\ihead{\footnotesize MAT 320 HW07}
\ohead{\footnotesize October 23, 2024}

\title{\vspace{-2em}MAT 320 HW07}
\author{\large Aditya Pahuja}
\date{\large Due 10/23}
\begin{document}
\maketitle

\section{Page of Lemmas}
\counterwithout{lemma}{problem}
\renewcommand{\thelemma}{L\arabic{lemma}}
The following two facts appeared frequently enough within the problems
and have long enough proofs
that I am moving them into this section for more convenient referencing.

\begin{lemma}[Open balls are open]
	Let $(X, d)$ be a metric space.
	For any $\eps > 0$, the open ball $B_X(x, \eps)$
	is actually an open set in $X$
	(thus open balls are worthy of their name).
\end{lemma}

\begin{proof}
	For any $y\in B_X(x,\eps)$, consider the open ball
	\[B_X(y, \eps - d_X(x, y))\]
	Any $z\in B_X(y, \eps - d_X(x,y))$ satisfies
	\[d_X(x, z) \le d_X(x, y) + d_X(y, z) < d_X(x, y) + \eps - d_X(x, y)
	= \eps.\]
	Thus, every element of $B_X(y,\eps-d_X(x,y))$ is in $B_X(x,\eps)$,
	meaning that, for every $y\in B_X(x,\eps)$,
	we have found an open ball $B_X(y,\eps-d_X(x,y))$ centered at $y$
	that is contained inside $B_X(x,\eps)$,
	proving that $B_X(x,\eps)$ is indeed open in $X$.
\end{proof}

\begin{lemma}[Linear combinations are continuous]
	Let $\mb x$ and $\mb y$ be arbitrary elements of $\RR^n$.
	Then, the function $\gamma\colon[0,1]\to\RR^n$ given by
	\[\gamma(t) = (1-t)\mb x + t\mb y\]
	is continuous (where $\RR^n$ has the standard metric,
	which we will call $d$).
\end{lemma}

\begin{proof}
	Let $t$ be in $[0,1]$; we need to show that $\gamma$ is continuous at $t$.
	Recall also that $\|\mb x\| = d(\mb x, \mb 0)$ is the norm of $\mb x$,
	so that $d(\mb x, \mb y) = \|\mb x - \mb y\|$,
	and that $\|c\mb x\| = |c|\cdot\|\mb x\|$ for any real $c$.
	
	Given $\eps > 0$, let $\delta = \frac{\eps}{d(\mb x,\mb y)} > 0$.
	If $|s - t| < \delta$, then we have
	\begin{align*}
		d(\gamma(s), \gamma(t)) &=
		\|((1-s)\mb x + s\mb y) - ((1-t)\mb x + t\mb y)\|\\
		&= \|(s-t)\mb y + (t-s)\mb x\|\\
		&= |s - t|\cdot \|\mb x - \mb y\|\\
		&= |s - t|\cdot d(\mb x, \mb y)\\
		&< \frac{\eps}{d(\mb x,\mb y)}\cdot d(\mb x,\mb y) = \eps.
	\end{align*}
	That is, we just found a $\delta > 0$ such that
	$|s - t| < \delta$ implies $d(\gamma(s),\gamma(t)) < \eps$,
	hence $\gamma$ is continuous at all $t\in[0,1]$.
\end{proof}

\begin{lemma}
	[Path connectedness is an equivalence relation]
	Given a metric space $(X, d)$, we say that
	$p\sim q$ if $p$ and $q$ are connected by a path,
	i.e. there exists a continuous function
	$\gamma\colon[0,1]\to X$ such that
	$\gamma(0) = p$ and $\gamma(1) = q$.
	Then, $\sim$ is an equivalence relation on $X$.
\end{lemma}

\begin{proof}
	First, the relation is reflexive.
	To construct a path from $a$ to itself,
	we can take $\gamma(t) = a$ for all $t\in[0,1]$,
	which is obviously continuous.
	
	Second, the relation is symmetric.
	If $a\sim b$, then there exists a path $\gamma$ going from $a$ to $b$.
	Noting that the map $t\mapsto 1-t$ is continuous, we have
	\[\widehat\gamma(t) = \gamma(1-t)\]
	is a continuous function $\widehat\gamma\colon[0,1]\to X$
	because it is a composition of continuous functions.
	It also satisfies $\widehat\gamma(0) = \gamma(1) = b$ and
	$\widehat\gamma(1) = \gamma(0) = a$, so $b\sim a$.
	
	Third, the relation is transitive.
	If $a\sim b$ and $b\sim c$, then
	there exist paths $\gamma_1$ and $\gamma_2$
	from $a$ to $b$ and from $b$ to $c$ respectively.
	To get a path from $a$ to $c$, we can ``glue'' the paths together:
	define the function $\gamma\colon[0,1]\to X$ given by
	\[\gamma(t) = \begin{cases}
		\gamma_1(2t) & \text{if }0 \le t \le \half\\
		\gamma_2(2t - 1) & \text{if }\half < t \le 1.
	\end{cases}\]
	Clearly $\gamma$ satisfies $\gamma(0) = \gamma_1(0) = a$
	and $\gamma(1) = \gamma_2(1) = c$.
	To show that $\gamma$ is continuous we consider three cases.
	\begin{itemize}
		\ii Let $0 \le t < \half$ and $\eps > 0$.
		Knowing that $\gamma_1$ is continuous at $2t$,
		there exists $0 < \delta < 1 - 2t$ such that
		\[|2x - 2t| < \delta \implies d(\gamma_1(2x), \gamma_1(2t)),\]
		where $\delta < 1 - 2t$ guarantees that
		$|2x - 2t| < \delta$ implies $2x < 1$.
		This means that $\gamma(t)$ is continuous at $t$, since now we have
		\[|x - t| < \frac{\delta}{2} \implies
		|2x - 2t| < \delta \implies d(\gamma_1(2x), \gamma_1(2t))
		= d(\gamma(x), \gamma(t)) < \eps,\]
		where we can write $\gamma(x)$ in place of $\gamma_1(2x)$
		because $2x < 1 \iff x < \half$ is forced.
		
		\ii Similarly, if $\half < t \le 1$, then, given $\eps > 0$,
		we pick $0 < \delta < 2t - 1$ such that
		\[|(2x - 1) - (2t - 1)| < \delta \implies
		d(\gamma_2(2x-1), \gamma_2(2t-1)) < \eps\]
		by continuity of $\gamma_2$. Here, $\delta < 2t-1$ guarantees that
		$|(2x-1) - (2t-1)| < \delta$ implies that $2x-1 > 0$.
		This translates to the continuity of $\gamma$ at $t$ because
		it says that
		\begin{align*}
			|x - t| < \frac{\delta}{2} &\implies
			|(2x-1) - (2t-1)| < \delta\\
			&\implies
			d(\gamma_2(2x-1), \gamma_2(2t-1)) =
			d(\gamma(x),\gamma(t)) < \eps,
		\end{align*}
		where we can safely substitute $\gamma_2(2x-1) = \gamma(x)$
		because $2x-1 > 0\iff x > \half$.
		
		\ii Finally, we need to check that $\gamma$ is continuous
		at $\half$. We know that $\gamma_1$ is continuous at $1$
		and $\gamma_2$ is continuous at $0$. Thus, we can say that,
		given $\eps > 0$, there exists $\delta_1 > 0$ such that
		\[|2x - 1| < \delta_1 \implies d(\gamma_1(2x), b) < \eps\]
		and there exists $\delta_2 > 0$ such that
		\[|(2x-1) - 0| < \delta_2 \implies d(\gamma_2(2x-1), b) < \eps,\]
		noting that $\gamma(\half) = \gamma_1(1) = \gamma_2(0) = b$.
		Therefore, if $\delta = \min(\delta_1, \delta_2)$, then
		\[\abs{x-\half} < \frac\delta2 \implies
		\abs{2x-1} < \delta \implies d(\gamma(x), b) < \eps\]
		where the second implication comes from
		the above two inequalities saying that
		if $x \le\half$ then
		$d(\gamma_1(2x), b) = d(\gamma(x),b)<\eps$
		and if $x > \half$ then
		$d(\gamma_2(2x-1),b) = d(\gamma(x),b)<\eps$.
		This means that $\gamma$ is continuous at $\half$, as desired.
	\end{itemize}
	Thus, $a$ and $c$ are connected by a path, which means $a\sim c$.
\end{proof}

\pagebreak
\section{Solutions to problems}
\counterwithin{lemma}{problem}

\begin{problem}
	Let $(X, d)$ be a metric space. Given $E\subseteq X$
	define the \vocab{closure} of $E$ as
	\[\ol E = \{x\in X : x\text{ is a limit point of $E$}\}\subseteq X.\]
	Explain why $E\subseteq \ol E$, and show that $\ol E$ equals
	the intersection of all closed subsets $F\subseteq X$
	that contain $E$.
\end{problem}

\begin{soln}
	First, to show that $E\subseteq \ol E$, we need to show that
	every point of $E$ is a limit point of $E$.
	This is obvious, though: for a given $x\in E$,
	we can find a sequence converging to $x$
	by taking the constant sequence $(x, x, x, \dots)$.
	
	Next, let $E'$ be the intersection of
	all closed subsets $F\subseteq X$ containing $E$;
	we need to show that $E' = \ol E$.
	Note that $E'$ is closed because it is the intersection of
	a collection of closed sets.
	Since $E'$ contains $E$, $E'$ must also contain every limit point of $E$
	because a set is closed if and only if it contains all of its limit points,
	which lets us conclude that $\ol E\subseteq E'$.
	
	On the other hand, we wish to show that $E'\subseteq \ol E$ as well.
	\begin{claim}
		$\ol E$ is a closed subset of $X$.
	\end{claim}
	\begin{proof}
		Let $(x_n)$ be a sequence of points in $\ol E$
		that converge to some $\ell\in X$,
		i.e. $\ell$ is a limit point of $\ol E$.
		For each $n\in\NN$, since $x_n$ is a limit point of $E$,
		we can find a sequence of points in $E$ that converge to $x_n$;
		in particular, we can find a point $y_n\in E$ such that
		$d(x_n, y_n) < \frac 1n$.
		By the triangle inequality,
		\[d(\ell, y_n) \le d(\ell, x_n) + d(x_n, y_n) <
		d(\ell, x_n) + \frac 1n\]
		Let $\eps > 0$ be an arbitrary positive real
		and let $N$ be a positive integer greater than $\frac 1\eps$.
		Then, $\eps - \frac 1N$ is a positive real number,
		hence $d(\ell, x_n) < \eps - \frac 1N$ is eventually true,
		say when $n > M$. This means
		\[d(\ell, x_n) + \frac 1n < \eps - \frac 1N + \frac 1n\]
		is true when $n > M$, and the expression on the right-hand side
		is smaller than $\eps$ when $n > N$, so the inequality
		\[d(\ell, y_n) < d(\ell, x_n) + \frac 1n <
		\eps - \frac 1N + \frac 1n < \eps\]
		holds eventually, namely when $n > \max(M, N)$.
		In particular, $(y_n)$ converges to $\ell$,
		so $\ell$ is a limit point of $E$ (because $y_n\in E$ for all $n$).
		This means $\ol E$ contains all of its limit points,
		so it is closed.
	\end{proof}
	This means $\ol E\supset E$ is a closed set containing $E$,
	which means $E'$ must be a subset of $\ol E$ because $E'$
	is a subset of every closed set containing $E$ by definition.
	This lets us conclude that $E' = \ol E$ as desired.
\end{soln}

\begin{problem}
	[Writing exercise]
	Show that a function between metric spaces $X\to Y$ is continuous
	if and only if $f\inv (U)$ is \emph{open}
	for each \emph{open} set $U\subseteq Y$.
	Show that the same is true if in this statement
	you replace the word \emph{open} with the word \emph{closed}.
\end{problem}

\begin{soln}
	Here's a picture:
	\begin{center}
		\begin{asy}
			unitsize(1.15cm);
			draw((1.2,0)..(0,2)..(-1.2,0)..(0,-2)..cycle);
			draw((6.2,0)..(5,2)..(3.8,0)..(5,-2)..cycle);
			pair x = (0.3,0.5);
			filldraw((0.3,0.08)..(0.6,0.2)..(0.72,0.37)..(0.3,1.5)..(-0.6,1.3)..(-0.4,0.2)..cycle,green+opacity(0.2),dashed);
			filldraw(circle(x,0.35),red+opacity(0.3),dashed);
			dot("$x$", x, dir(90));
			pair fx = (4.8, -.3);
			filldraw(circle(fx,0.7),green+opacity(0.2),dashed);
			filldraw((5.1,-.34)..(4.8,-0.1)..(4.5,0.2)..(4.4,-0.2)..(5.2,-0.7)..cycle,red+opacity(0.3),dashed);
			dot("$f(x)$", fx, dir(-70),fontsize(9));
			draw((0.7,0.5)..(2.5,0.3)..(4.3,-0.2),EndArrow(5));
			
			label("$B_X(x,\delta)$",x+(0,0.3),dir(100),fontsize(7.5));
			label("$B_Y(f(x),\varepsilon)$",fx-(0,0.7),dir(-80),fontsize(7.5));
			label("$f^{-1}(U)$",(0,2),dir(90),fontsize(11));
			label("$U$",(5,2),dir(90),fontsize(11));
			label("$f$",(2.5,0.3),dir(90));
			label("$f^{-1}$",(2.5,0.15),dir(-90));
			label("$f(B_X(x,\delta))$", fx+(0.6,0.4), fontsize(7.5));
			draw((0.8,0.35)..(2.5,0.15)..(4.05,-0.28),BeginArrow(5));
			label("$f^{-1}(B_Y(f(x),\varepsilon))$",(0,0.1),dir(-90),fontsize(7.5));
		\end{asy}
	\end{center}
	
	First, suppose $f$ is continuous.
	For an open $U\subseteq Y$, let $x\in f\inv(U)$
	and let $B_Y(f(x),\eps)$ be contained in $U$.
	Then, there exists a $\delta > 0$ so that
	\[f(B_X(x,\delta))\subseteq B_Y(f(x),\eps)\]
	since being continuous at $x$ implies the existence of $\delta> 0$
	such that $d_X(x, y) < \delta$ implies $d_Y(f(x), f(y)) < \eps$.
	This means that $B_X(x,\delta)\subseteq f\inv(U)$, i.e.
	for each $x\in f\inv(U)$, we can find an open ball centered at $x$
	contained in $U$, so $f\inv(U)$ is open.
	
	Conversely, suppose $f\inv(U)$ is open in $X$ for all open $U\subseteq Y$.
	By Lemma L1, the open ball $B_Y(f(x),\eps)$ is open in $Y$.
	Thus, $f\inv(B_Y(f(x),\eps))$ is open in $X$ and contains $x$.
	This means we can find an open ball
	$B_X(x,\delta)\subseteq f\inv(B_Y(f(x),\eps))$ due to openness,
	so, for each $\eps>0$, there exists $\delta>0$ such that
	\[f(B_X(x,\delta))\subseteq B_Y(f(x),\eps),\]
	so $f$ is continuous.\\
	
	Now, we prove the second part (replacing ``open'' with ``closed'').
	The statement
	\[E\subseteq Y\text{ is closed}\implies
	f\inv(E)\subseteq X\text{ is closed}\]
	is equivalent to
	\[Y\setdiff E\subseteq Y\text{ is open} \implies
	f\inv(Y\setdiff E) = X\setdiff f\inv(E)\subseteq X\text { is open}.\]
	Then, $Y\setdiff E$ is open iff
	$E = Y\setdiff U$ for an open $U\subseteq Y$,
	so the line above is equivalent to
	\[U\subseteq Y\text{ is open}\implies f\inv(U)\subseteq X\text{ is open}.\]
	This is then equivalent to $f$ being continuous
	by the work done above, so we are done.
\end{soln}

\begin{problem}
	Show that the composition of two continuous functions
	$f\colon X\to Y$ and $g\colon Y\to Z$ between metric spaces
	gives rise to a continuous function $g\circ f\colon X\to Z$.
\end{problem}

\begin{soln}
	Let $V$ be an open subset of $Z$. Then, letting $U = g\inv(V)$,
	\[(g\circ f)\inv(V) = f\inv(g\inv(V)) = f\inv(U)\subseteq X.\]
	By Problem 2, since $g$ and $f$ are continuous,
	$U = g\inv(V)$ is open in $Y$ and $f\inv(U)$ is open in $X$,
	so, by Problem 2 again, $f\circ g$ is continuous.
\end{soln}

\begin{problem}
	Let $(X, d_X)$ be a metric space.
	Show that a subset $A\subseteq E$ is open in the induced metric
	$d_E$ if and only if there exists an open set $U\subseteq X$
	such that $A = U\cap E$.
	Show that if $E$ is open in $X$ then the open subsets
	$(E, d_E)$ are exactly the open subsets of $X$ that are contained in $E$.
	Show that there exist open subsets of $[0, +\infty)$
	that are not open in $\RR$.
\end{problem}

\begin{soln}
	We first show that $A\subseteq E$ is open if and only if $A = U\cap E$
	where $U$ is open in $X$.
	
	If $A\subseteq E$ is open with respect to the induced metric,
	then, for each $x\in A$, there exists an $\eps_x > 0$
	such that the open ball
	\[B_E(x, \eps_x) = \{y\in E : d_E(x, y) < \eps_x\}\]
	is a subset of $A$.
	We can therefore write
	\[A = \bigcup_{x\in A} B_E(x, \eps_x),\]
	as, on one hand, each of the balls is contained in $A$,
	thus so is their union, so the right-hand side is a subset of $A$;
	and, on the other hand, each $x\in A$ is in the open ball $B_E(x, \eps_x)$,
	so the union contains $A$.
	
	However, $B_E(x, \eps_x)$ can be written as $B_X(x, \eps_x) \cap E$,
	noting that $d_E(x, y) = d_X(x, y)$ for $x,\,y\in E$.
	We can thus write the above as
	\[A = \bigcup_{x\in A} (B_X(x, \eps_x)\cap E) =
	E\cap \left(\bigcup_{x\in A} B_X(x, \eps_x)\right).\]
	Thus, since each open ball is open in $X$ by Lemma L1,
	the union of the open balls is an open set $U$,
	and so we have found an open set $U$ such that $A = E\cap U$.
	
	Conversely, if $U\subseteq X$ is open, then for each $x\in U$,
	there exists $\eps_x > 0$ such that the ball $B_X(x, \eps_x)$
	is a subset of $U$. This means that, if $x\in U\cap E$,
	\[B_E(x, \eps_x) = B_X(x, \eps_x)\cap E
	 \subseteq U\cap E.\]
	That is, for each $x$ in the intersection of $U$ and $E$,
	we can find an open ball in the induced metric $d_E$
	centered at $x$ that is contained in $U\cap E$,
	which is what it means for $U\cap E$ to be open in $E$.\\
	
	If $E$ is open in $X$, then $A\subseteq E$ is open in $E$
	if and only if $A = E\cap U$ for some $U$ that is open in $X$.
	However, this means that $A$, being the intersection of
	two open subsets of $X$, is also open in $X$,
	which means $A$ is open in $E$ if and only if $A$ is open in $X$.\\
	
	Finally, the set $U = [0, 1)$ is not open in $\RR$ because
	any open ball centered at $0$ must contain numbers outside $U$,
	but it is an open subset of $[0, +\infty)$ because
	$U$ is the intersection of the open set $(-1, 1)\subseteq\RR$
	with $[0, +\infty)$.
\end{soln}

\begin{problem}
	Explain the following statement that was made during class:
	$E\subseteq X$ is disconnected if and only if
	there are open sets $A,\, B\subseteq X$ such that
	$E\subseteq A\cup B$, $E\cap A\ne\varnothing$, $E\cap B\ne\varnothing$,
	and $E\cap(A\cap B) = \varnothing$.
	How is this statement simplified if $E$ is open?
\end{problem}

\begin{soln}
	By definition, $E\subseteq X$ is disconnected if and only if
	there exist nonempty disjoint open subsets $A'$ and $B'$ of $E$
	whose union is $E$. By Problem 5, we know that
	$A' = A\cap E$ and $B' = B\cap E$ where $A$ and $B$ are open in $X$.
	We can now verify the criteria for $E$ being disconnected:
	\begin{itemize}
		\ii The condition that $A'$ and $B'$ are both nonempty
		translates to $A\cap E$ and $B\cap E$ both being nonempty.
		\ii The condition that $A'$ and $B'$ are disjoint translates to
		$(A\cap E)\cap(B\cap E) = (A\cap B)\cap E$ being the empty set.
		\ii Finally, $A'\cap B' = E$ translates to $(A\cap E)\cup(B\cap E) =
		(A\cup B)\cap E = E$, which is true if and only if $E\subseteq A\cup B$.
	\end{itemize}
	
	If $E$ is open, then Problem 5 says that
	$A'$ and $B'$ will also be open in $X$,
	so the condition for disconnectedness is just that
	there exist disjoint open subsets of $X$ whose union is $E$.
\end{soln}

\begin{problem}
	Let $(X, d)$ be a metric space. Show that if $D\subseteq X$ is connected
	then its closure is also connected.
	Use this to show that, given $(a, b)$ is connected,
	$(a, b]$, $[a, b)$, and $[a, b]$ are also connected.
\end{problem}

\begin{soln}
	We prove the contrapositive statement that
	$\ol D$ being disconnected implies $D$ is disconnected.
	If $\ol D$ is disconnected, then, by Problem 5, there exist
	open subsets of $X$, say $A$ and $B$, whose union contains $\ol D$
	such that $A\cap\ol D$ and $B\cap\ol D$ are both nonempty
	but $A\cap B\cap\ol D$ is empty.
	Since $D\subseteq \ol D$ by Problem 1, we already know that
	\[A\cap B\cap D = \varnothing,\qquad D\subseteq A\cap B.\]
	To show that $D$ is disconnected we need to prove that
	$A\cap D$ and $B\cap D$ are nonempty.
	We will prove that $A\cap D$ is nonempty;
	the proof for $B\cap D$ is identical.
	
	Suppose that $x$ is an element of $A\cap\ol D$.
	Since $A\cap\ol D$ is open in $\ol D$,
	there exists an $\eps > 0$ such that the open ball $B_{\ol D}(x, \eps)$
	is a subset of $A\cap\ol D$. Then, $x$ is a limit point of $D$,
	so we can find a sequence of points in $D$ converging to $x$.
	In particular, we can find a $y\in D$ such that $y\in B_{\ol D}(x,\eps)$,
	which means that $y$ is in $A\cap\ol D$. This means that $y\in A$,
	so $y\in A\cap D$, and thus $A\cap D$ is nonempty, as desired.
	
	Now, we are ready to show that all kinds of intervals are connected!
	First, for closed intervals, we show that $[a, b]$ is the closure of
	$(a, b)$ over $\RR$. Indeed, $a$ and $b$ are limit points of $(a, b)$
	because the sequences
	\[a_n = a + \frac{b-a}n,\qquad b_n = b - \frac {b-a}n\]
	for $n > 1$ converge to $a$ and $b$ respectively
	and consist only of points in $(a, b)$
	(since $a < a_n < a + b - a = b$ and $a = b - (b - a) < b_n < b$),
	and obviously all points of $(a, b)$ are in its closure by Problem 1.
	Moreover, $(a, b)$ has no other limit points because,
	if a sequence $(x_n)$ satisfies $a < x_n < b$ for all $n$, then
	taking limits as $n$ goes to infinity shows that
	\[a \le \lim_{n\to\infty}x_n \le b\]
	if the limit exists. Thus $[a, b] = \ol{(a, b)}$ is connected.
	
	For half-open intervals, we use a little trick.
	\begin{lemma}
		Consider an open set $U\subseteq X$.
		A subset $E\subseteq U$ is connected in the metric space $(U, d_U)$
		if and only if it is connected in $(X, d)$.
	\end{lemma}
	\begin{proof}
		First, suppose $E$ is disconnected in $(X, d)$.
		Then, there exist open subsets of $X$, say $A$ and $B$, such that
		\[E\subseteq A\cup B,\quad
		E\cap A\ne\varnothing, \quad E\cap B\ne\varnothing,\quad
		E\cap(A\cap B)=\varnothing.\]
		Then, $A' = A\cap U$ and $B' = B\cap U$ are open in $U$,
		and, since $E\subseteq U$, we can say that
		\begin{itemize}
			\ii $E$ is a subset of $A'\cup B' = (A\cup B)\cap U$,
			\ii $E\cap A' = E\cap A\ne \varnothing$,
			\ii $E\cap B' = E\cap B\ne\varnothing$, and
			\ii $E\cap(A'\cap B')\subseteq E\cap(A\cap B) = \varnothing$
			(which obviously means $E\cap(A'\cap B')=\varnothing$).
		\end{itemize}
		Thus, $E$ is also disconnected in $(U, d_U)$.
		
		For the other direction, suppose $E$ is disconnected in $(U, d_U)$.
		Then, $E$ is a union of disjoint sets $A$ and $B$ that are open in $U$.
		Problem 4 shows us that, since $U$ is open, $A$ and $B$ are open in $X$,
		so we just treat $A$ and $B$ as open subsets of $X$ instead of $U$
		to say that $E$ is connected in $X$.
	\end{proof}

	Now consider the subspaces of $\RR$ given by
	$A = (a, +\infty)\subseteq\RR$ and $B = (-\infty, b)\subseteq\RR$,
	which are open intervals and thus open.
	Clearly, $(a, b)$ is still connected within each of these spaces.
	The closure of $(a, b)$ within $A$ is now $[a, b]\cap A = (a, b]$,
	since $a\notin A$ (but all the other limit points are exactly preserved),
	and similarly the closure of $(a, b)$ within $B$ is $[a, b)$.
	The lemma tells us that the connectedness of these half-open intervals
	within the respective subspace
	is equivalent to their connectedness in $\RR$ because $A$ and $B$ are open,
	so $(a, b]$ and $[a, b)$ are connected in $\RR$, too!
\end{soln}

\begin{problem}
	Let $f\colon X\to Y$ be a continuous function between metric spaces.
	Show that if $E\subseteq X$ is path connected,
	then $f(E)\subseteq Y$ is also path connected.
\end{problem}

\begin{proof}
	Given two points $a$ and $b$ in $f(E)$,
	there exist $a'$ and $b'$ in $E$ such that
	$f(a') = a$ and $f(b') = b$.
	Let $\gamma\colon[0,1]\to E$ be a path between $a'$ and $b'$,
	so $\gamma(0) = a'$, $\gamma(1) = b'$, and $\gamma$ is continuous.
	Then, define $\Gamma\colon[0,1]\to f(E)$ by
	$\Gamma(t) = f(\gamma(t))$. This is a composition of continuous functions,
	so it is continuous by Problem 3, and
	\[\Gamma(0) = f(\gamma(0)) = f(a') = a,\qquad
	\Gamma(1) = f(\gamma(1)) = f(b') = b.\]
	Thus, we have constructed a path $\Gamma$ between $a$ and $b$,
	so any pair of points in $f(E)$ has a path between them,
	and thus $f(E)$ is path connected.
\end{proof}

\begin{problem}
	A subset $C\subseteq \RR^n$ is called \vocab{convex} if
	the line segment connecting any two points of $C$
	lies completely within $C$.
	Show that any convex subset of $\RR^n$ is connected.
	Is it true that a subset of $\RR^n$ is connected if and only if
	it is convex?
\end{problem}

\begin{soln}
	Let $C\subseteq \RR^n$ be convex and let $\mb x,\,\mb y\in C$.
	The line segment between $\mb x$ and $\mb y$ is the set
	\[\ell(\mb x, \mb y) = \{(1 - t)\mb x + t\mb y : t\in[0, 1]\}\subseteq C.\]
	Define $\gamma\colon[0, 1]\to C$ by
	\[\gamma(t) = (1 - t)\mb x + t\mb y.\]
	This is a continuous function by Lemma L2 such that
	$\gamma(0) = \mb x$ and $\gamma(1) = \mb y$,
	so it is a path between $\mb x$ and $\mb y$.
	Thus, we can find a path between any two points of $C$,
	so $C$ is path connected, meaning it is connected.
	
	The converse of this statement isn't true ---
	there exist connected sets in $\RR^n$ that aren't convex.
	For example, in $\RR^2$, we will show that
	\[S = \{(a, 0) : a\in[0, 1]\}\cup \{(0, a) : a\in[0, 1]\}\]
	is connected but not convex.
	\begin{center}
		\begin{asy}
			draw((-1,0)--(3,0),linewidth(0.25),Arrows(5));
			draw((0,-1)--(0,3),linewidth(0.25),Arrows(5));
			draw((0,2)--(0,0)--(2,0), linewidth(1.2)+red);
			draw((0,2)--(2,0),dashed);
			dot((0,2),red);
			dot((2,0),red);
		\end{asy}
	\end{center}

	Clearly $S$ isn't convex, because the segment between $(0,1)$
	and $(1,0)$ contains points outside of $S$, like $(\half,\half)$.
	
	However, we will show that it is path connected and thus connected.
	Let $\mb a$ and $\mb b$ be two points in $S$.
	\begin{itemize}
		\ii First, we will show that
		$\mb 0$ is connected to every point $\mb x\in S$.
		Define $\gamma\colon[0,1]\to S$ by $\gamma(t) = t\mb x$.
		If $\mb x\in S$ then $t\mb x\in S$ because
		$\mb x = (a, 0)$ or $(0, a)$ for some $a\in[0, 1]$,
		and $t\mb x = (ta, 0)$ or $(0, ta)$, in which case
		$ta\in[0, 1]$ because $0\le t\le1$.
		Additionally, writing $\gamma(t) = (1-t)\mb 0 + t\mb x$,
		we can use Lemma L2 to conclude that $\gamma$ is continuous,
		and we can obviously see that $\gamma(0) = \mb 0$ and
		$\gamma(1) = \mb x$.
		We have thus verified all the conditions for $\gamma$
		to be a path between the two points,
		so $\mb 0$ and $\mb x$ are connected by a path.
		
		\ii Now, $\mb a$ and $\mb b$ are connected by a path to $(0,0)$,
		so Lemma L3 lets us say that, since
		path connectedness is transitive, $\mb a$ and $\mb b$
		are also connected by some path.
	\end{itemize}
	Thus $S$ is (path) connected, but not convex.
\end{soln}

\begin{problem}
	Show that an open set $E\subseteq \RR^n$ is connected
	if and only if it is path connected.
\end{problem}

\begin{soln}
	We have already shown in class that $E$ being path connected
	implies that it is connected.
	We wish to prove the converse statement that
	if $E$ isn't path connected, then it is not connected.

	\begin{claim}
		Let $\eps > 0$ and $\mb x\in\RR^n$.
		If $\mb y \in B_{\RR^n}(\mb x, \eps)$,
		then there exists a path connecting $\mb x$ and $\mb y$.
	\end{claim}
	\begin{proof}
		Noting that we are working in $\RR^n$,
		we consider the line segment connecting $\mb x$ and $\mb y$,
		consisting of the points
		$\gamma(t) = (1-t)\mb x+t\mb y$ with $0\le t\le 1$.
		Observe that
		\begin{align*}
			d_{\RR^n}(\mb x,\gamma(t)) &= d(\mb x, (1-t)\mb x + t\mb y)\\
			&= \|(1 - t)\mb x +t\mb y - x\| = \|t(\mb y - \mb x)\|\\
			&= t\cdot d_{\RR^n}(\mb x,\mb y) < t\eps \le \eps,
		\end{align*}
		noting that $t\eps \le \eps$ because $0\le t\le1$.
		Thus, the line segment between $\mb x$ and $\mb y$
		is always contained in the open ball,
		so we can just take this segment as a path between the two points,
		noting that Lemma L2 tells us that $\gamma(t)$ is continuous.
	\end{proof}
	
	As in Lemma L3, define the relation $\sim$ on $E$ by saying that
	$a\sim b$ if and only if there is a path inside $E$
	that connects $a$ to $b$.
	Let $\mathcal U$ be the collection of
	equivalence classes of $E$ under $\sim$;
	the elements of $\mathcal U$ form a partition of $E$.
	We must have more than one element in $\mathcal U$
	because if there is only one equivalence class under $\sim$,
	then any two elements of $E$ are connected by a path,
	i.e. $E$ is path connected, which contradicts our assumption.
	
	\begin{claim}
		Each element of $\mathcal U$ is an open subset of $E$.
	\end{claim}
	\begin{proof}
		Let $U$ be an element of $\mathcal U$ and let $x$ be in $U$.
		We know that, because $E$ is open in $\RR^n$,
		there exists $\eps > 0$ such that
		the open ball $B_{\RR^n}(x, \eps)$ is contained in $E$.
		By Claim 9.1, each element of this ball
		can be connected to $x$ via a path,
		so each element of this ball is in $U$, meaning
		\[B_{\RR^n}(x, \eps) \subseteq U.\]
		This holds for every $x\in U$, so $U$ is open in $\RR^n$
		and thus also open in $E$
		(because $U\subseteq E$ implies $U = E\cap U$),
		which is what we wanted to show.
	\end{proof}

	To finish, $\mathcal U$ is a partition of $E$
	into (more than one) open sets --- that is, $E$ is the union
	of a collection of pairwise disjoint open sets.
	Let $U_1$ be an arbitrary element of $\mathcal U$
	(which is obviously nonempty)
	and let $U_2$ be the union of all the other elements of $\mathcal U$, so
	\[U_1\cup U_2 = \bigcup_{U\in\mathcal U} U = E.\]
	Because $\cal U$ has more than one element, $U_2$ is nonempty.
	Additionally, $U_2$ is a union of open sets and thus also open in $E$,
	and it is disjoint from $U_1$ because the elements of $\mathcal U$
	are pairwise disjoint.
	Therefore, we have found sets $U_1$ and $U_2$ that are open in $E$,
	nonempty and disjoint, and whose union is $E$,
	which proves that $E$ is disconnected.
\end{soln}

\begin{problem}
	Show that if $X$ and $Y$ are path connected metric spaces
	then so is $X\times Y$.
	Extra: do the same but replacing
	\emph{path connected} with \emph{connected}.
\end{problem}

\begin{soln}
	First, suppose $X$ and $Y$ are path connected, and let
	$(a_1,b_1)$ and $(a_2, b_2)$ be in $X\times Y$.
	We know that there exist paths
	$\gamma_X\colon[0, 1]\to X$ connecting $a_1$ and $a_2$, and
	$\gamma_Y\colon[0,1]\to Y$ connecting $b_1$ and $b_2$.
	Then, I claim that $\gamma\colon[0, 1]\to X\times Y$,
	given by
	\[\gamma(t) = (\gamma_X(t), \gamma_Y(t))\]
	is a path between $(a_1,b_1)$ and $(a_2,b_2)$.
	Obviously $\gamma(0) = (a_1, b_1)$ and $\gamma(1) = (a_2, b_2)$,
	so we only need to prove that $\gamma$ is continuous
	with respect to the product metric $d_{X\times Y}$
	(which we can recall is defined as
	$d_{X\times Y}((a_1,b_1),(a_2,b_2)) = d_X(a_1, a_2) + d_Y(b_1, b_2)$).
	The continuity of $\gamma_X$ and $\gamma_Y$ amount to the following:
	for each $t\in[0,1]$ and $\eps>0$, there exists $\delta_1>0$ such that
	\[|x - t| < \delta_1\implies d_X(\gamma_X(x), \gamma_X(t)) < \frac\eps2,\]
	and there exists $\delta_2>0$ such that
	\[|y - t| < \delta_2\implies d_Y(\gamma_Y(y), \gamma_Y(t)) < \frac\eps2.\]
	Then, if $\delta = \min(\delta_1,\delta_2)$, we find that
	$|s - t| < \delta$ implies that
	\begin{align*}
		d_{X\times Y}(\gamma(s), \gamma(t)) &=
		d_{X\times Y}((\gamma_X(s), \gamma_Y(s)), (\gamma_X(t),\gamma_Y(t)))\\
		&= d_X(\gamma_X(s),\gamma_X(t)) + d_Y(\gamma_Y(s),\gamma_Y(t))\\
		&< \frac\eps2 + \frac\eps2 = \eps,
	\end{align*}
	hence $\gamma$ is continuous at every $t\in[0,1]$, so it is continuous.
	This proves that $X$ and $Y$ being path connected
	implies that $X\times Y$ is path connected.\\
	
	Now we show that if $X$ and $Y$ are connected, then so is $X\times Y$.

	\begin{lemma}
		Let $U$ be a subset of $X\times Y$.
		For a given $y\in Y$,
		define $\pi_X(U, y)$ to be the set $\{x : (x, y)\in U\}\subseteq X$, and
		for a given $x\in X$,
		define $\pi_Y(U, x)$ to be the set $\{y : (x, y)\in U\}\subseteq Y$.
		Then, if $U$ is open in $X\times Y$,
		then, for any $x\in X$ and $y\in Y$,
		$\pi_X(U, y)$ and $\pi_Y(U, x)$ are open
		in $X$ and $Y$ respectively.
	\end{lemma}
	\begin{proof}
		Suppose $U$ is open in $X\times Y$.
		Then, for some $y\in Y$, let $L_y = \{(x, y) : x\in X\}$.
		Observe that the induced metric on $L_y$ is
		\[d_{L_y}((a, y), (b, y)) = d_{X\times Y}((a,y), (b,y)) =
		d_X(a, b) + d_Y(y, y) = d_X(a, b).\]
		We see that $L_y\cap U$ is the set
		\[\{(x, y) : x\in \pi_X(U, y)\} = \pi_X(U, y)\times\{y\},\]
		which is open in $L_y$ with the induced metric.
		In particular, for any $a\in \pi_X(U, y)$,
		there exists an $\eps$ such that
		$B_{L_y}((a,y), \eps)\subseteq L_y\cap U$. That is,
		\[d_{L_y}((a, y), (x, y)) = d_X(a, x) < \eps \implies
		(x, y)\in L_y\cap U = \pi_X(U, y)\times\{y\},\]
		but this is the same thing as saying that $x\in \pi_X(U, y)$.
		To be clear, we just showed that, for each $a\in\pi_X(U, y)$,
		there exists an $\eps$ such that
		$d_X(a, x) < \eps$ implies $x\in \pi_X(U, y)$,
		so $B_X(a, \eps)\subseteq \pi_X(U, y)$.
		Thus, $\pi_X(U, y)$ is open in $X$.
		
		The proof that $\pi_Y(U, x)$ is open in $Y$ for any $x\in X$
		is identical, so we are done.
	\end{proof}

	This lemma is our mechanism for partitioning $X$ and $Y$ into open sets.
	Specifically, suppose $X\times Y$ is disconnected,
	so that it can be written as a union of nonempty disjoint sets $A$ and $B$
	that are open in $X\times Y$;
	we want to show that one of $X$, $Y$ must be disconnected.
	The lemma says that for any $x\in X$ and $y\in Y$, the four sets
	\[\pi_X(A, y),\quad \pi_X(B, y),\quad\pi_X(A, x),\quad \pi_X(B, x)\]
	are all open.
	Note also that $\pi_X(A, y)$ is the complement of $\pi_X(B, y)$ in $X$ 
	because they are obviously disjoint (since $A$ and $B$ are disjoint) and
	\[\pi_X(A, y)\cup \pi_X(B, y) = \{x : (x, y)\in A\cup B\}
	= \{x : (x, y)\in X\times Y\} = X.\]
	Similarly, $\pi_Y(A, x)$ is the complement of $\pi_Y(B, x)$ in $Y$.
	In other words, $\pi_X$ and $\pi_Y$ naturally partition $X$ and $Y$
	into disjoint open sets.
	Thus, we just need to show that there either exists a $y$ such that
	$\pi_X(A, y)$ and its complement $\pi_X(B, y)$ are both nonempty, or
	there exists an $x$ such that
	$\pi_Y(A, x)$ and its complement $\pi_Y(B, x)$ are both nonempty.
	Equivalently,
	we want $\pi_X(A, y)$ to be both nonempty and not equal to $X$,
	or we want $\pi_Y(A, x)$ to be both nonempty and not equal to $Y$.
	
	Now, if we could find a $y$ that made $\pi_X(A, y)$ both nonempty
	and not equal to $X$ (so $X$ is disconnected),
	then we'd be done, so suppose no such $y$ exists.
	We will then prove that $Y$ is disconnected.
	Observe that
	\[\bigcup_{y\in Y} \pi_X(A,y)\times\{y\} =
	\bigcup_{y\in Y} \{(x,y) : (x, y)\in A\}
	= \{(x,y) : (x,y)\in A\} = A,\]
	so if $\pi_X(A,y) = \varnothing$ for all $y$ then $A$ would be empty,
	and if $\pi_X(A,y) = X$ for all $y$ then $A$ would equal $X\times Y$,
	both of which contradict our assumptions about $A$.
	Therefore there exist $y_1$ and $y_2$ in $Y$ such that
	$\pi_X(A, y_1) = \varnothing$ and $\pi_X(A, y_2) = X$,
	so we know that $(x, y_1)\notin A$ and $(x, y_2)\in A$ for any $x\in X$.
	Hence, for any choice of $x\in X$,
	\[\pi_Y(A, x) = \{y: (x,y)\in A\},\]
	so $y_1\notin\pi_Y(A, x)$ but $y_2\in\pi_Y(A, x)$.
	In particular $\pi_Y(A, x)$ is nonempty and not equal to $Y$
	(thus having nonempty complement), which proves that $Y$ is disconnected.
\end{soln}

\begin{problem}
	In $\ell^2$ consider the infinite-dimensional unit sphere:
	\[S^\infty = \left\{\mb x\in\ell^2 : \sum_{i=1}^\infty x_i^2 = 1\right\}.\]
	Is $S^\infty\subseteq\ell^2$ a connected subspace of $\ell^2$?
\end{problem}

\begin{soln}
	We show that $S^\infty$ is path connected.
	Let $\mb x$ and $\mb y$ be points in $S^\infty$
	such that $\mb x + \mb y\ne \mb 0$.
	Then, I claim that the following function $\gamma\colon[0,1]\to S^\infty$
	is a path between $\mb x$ and $\mb y$:
	\[\gamma(t) = \frac{(1-t)\mb x + t\mb y}{\|(1-t)\mb x + t\mb y\|}.\]
	(The norm $\|\mb z\|$ is $d(\mb z, \mb 0)
	= \sqrt{\sum_{n=1}^\infty z_i^2}$ and has
	roughly the same properties as the $\RR^n$ norm, namely
	$d(\mb x,\mb y) = \|\mb x - \mb y\|$ and
	$\|c\mb x\| = |c|\cdot\|\mb x\|$ for any $c\in\RR$.)
	Some preliminary checks:
	\begin{itemize}
	 	\ii First,
	 	\[\gamma(0) = \frac{\mb x}{\|\mb x\|} = \mb x,\qquad
	 	\gamma(1) = \frac{\mb y}{\|\mb y\|} = \mb y\]
	 	because the norms of $\mb x$ and $\mb y$ are both $1$.
	 	
	 	\ii Second, $\gamma(t)$ is always defined. The only concern
	 	is that the denominator may be zero, in which case
	 	$(1-t)\mb x + t\mb y$ would be the zero vector (denoted $\mb 0$).
	 	If $t = 0$, then $\gamma(t) = \mb x\ne\mb 0$, and, for $t\ne 0$,
	 	\[(1-t)\mb x + t\mb y = \mb 0 \iff \mb x = \frac{1-t}{t}\mb y.\]
	 	We know that $\|\mb x\| = \|\mb y\| = 1$,
	 	so in order for the norms of $\mb x$ and $\frac{1-t}{t}\mb y$
	 	to be equal, we must have $\frac{1-t}{t} = 1$, so $t = \half$,
	 	which means that the denominator can only be zero if
	 	$\mb x + \mb y = 0$, and this is precisely the case I excluded.
	 	
	 	\ii Third, $\gamma(t)$ is always in $S^\infty$ because
	 	the norm of $\gamma(t)$ is just
	 	\[\frac{1}{\|(1-t)\mb x + t\mb y\|}\cdot\|(1-t)\mb x + t\mb y\| = 1.\]
	\end{itemize}
	
	To prove that $\gamma$ is indeed a path between $\mb x$ and $\mb y$,
	we have to show that $\gamma$ is continuous.
	We have shown in Lemma L2 that $(1-t)\mb x + t\mb y$
	is a continuous function of $t$.
	In fact the function $\mb z\mapsto\|\mb z\|$
	is also continuous because, for each $\eps > 0$,
	we can take $\delta = \eps$ to get
	\[d(\mb x, \mb y) < \eps \implies
	\abs{\|\mb x\| - \|\mb y\|} \le \|\mb x - \mb y\| = d(\mb x,\mb y) < \eps,\]
	where the inequality of norms follows from the triangle inequality:
	if WLOG $\|\mb x\| \ge \|\mb y\|$, then
	\[\|\mb x\| = d(\mb x,\mb 0) \le d(\mb x, \mb y) + d(\mb y, \mb 0)
	= \|\mb x - \mb y\| + \|\mb y\| \implies
	\|\mb x\| - \|\mb y\| \le \|\mb x-\mb y\|.\]
	We can then show that $\gamma$ is continuous as follows:
	\begin{itemize}
		\ii First, $\|(1-t)\mb x + t\mb y\|$ is the composed map
		$[0,1]\to \ell^2\to \RR$ given by
		\[t\mapsto (1-t)\mb x + t\mb y\mapsto \|(1-t)\mb x + t\mb y\|,\]
		which is continuous because the intermediate maps are both continuous,
		as per the result of Problem 3.
		\ii Then, since $\|(1-t)\mb x + t\mb y\|$ is never zero for $t\in[0,1]$,
		the quotient of continuous functions
		\[\frac{(1-t)\mb x + t\mb y}{\|(1-t)\mb x + t\mb y\|} = \gamma(t)\]
		is also a continuous function of $t$.
	\end{itemize}
	Therefore, we have shown that if $\mb x$ and $\mb y$ are not
	antipodal points on $S^\infty$, then they are connected by a path.
	
	To finish, in the case where $\mb x = -\mb y$,
	pick an arbitrary third point $\mb p$ that is not equal to
	$\mb x$ or $\mb y$. Then $\mb p$ is not equal to
	$-\mb x = \mb y$ or $-\mb y = \mb x$, so $\mb x$ and $\mb y$
	are both connected by a path to $\mb p$. By Lemma L3, this implies
	$\mb x$ and $\mb y$ are connected by a path
	as the relation of ``connected by a path'' is transitive.
	
	Therefore, any two points in $S^\infty$ are connected by a path,
	so $S^\infty$ is path connected and thus connected.
\end{soln}

\begin{problem}
	Explain why $\QQ$ is dense in $\RR$ with the standard metric.
\end{problem}

\begin{soln}
	We have shown that any real number
	is the limit of some sequence of rational numbers,
	so every real number is a limit point of $\QQ$;
	thus, $\ol\QQ = \RR$.
\end{soln}

\begin{problem}
	Show that
	\[D = \{\mb x\in\RR^\infty :
	x_i = 0\text{ for all but finitely many } i\in\NN\}\]
	is dense in $\ell^1$.
\end{problem}

\begin{soln}
	Let $\mb x$ be in $\ell^1$ and define $s = \sum_{k=1}^\infty |x_k|$.
	I will show that $\mb x$ is a limit point of $D$
	by proving that it is the limit of the sequence $(\mb x_n)$,
	where $\mb x_n$ is the sequence whose $k$th term is $x_k$
	when $k\le n$, and $0$ when $k > 0$
	(for example $\mb x_3 = (x_1, x_2, x_3, 0, 0, \dots)$).
	We see that
	\[d(\mb x, \mb x_n) = \sum_{k=1}^n |x_k - x_k| +
	\sum_{k=n+1}^\infty |x_k - 0| = \sum_{k=n+1}^\infty |x_k| =
	s - \sum_{k=1}^n |x_k| = s - s_n,\]
	where we define $s_n$ to be the partial sums of the sequence $(|x_k|)$.
	Taking the limit, we see that
	\[\lim_{n\to\infty} d(\mb x, \mb x_n) =
	\lim_{n\to\infty} s - s_n = s - \lim_{n\to\infty} s_n = s - s = 0,\]
	so $d(\mb x, \mb x_n)$ will eventually be smaller than $\eps$
	for any choice of $\eps > 0$; thus, $(\mb x_n)$ converges to $\mb x$
	in the metric space $(\ell^1, d)$.
\end{soln}

\begin{problem}
	In the last homework we proved that $\ell^1$ is a subset of $\ell^2$.
	Is $\ell^1$ dense in $\ell^2$?
\end{problem}

\begin{soln}
	In fact
	\[D = \{\mb x\in\RR^\infty :
	x_i = 0\text{ for all but finitely many } i\in\NN\}\subseteq \ell^1\]
	is dense in $\ell^2$, too --- this is enough to prove
	$\ell^1$ is dense in $\ell^2$ because a limit point of $D$
	is also a limit point of $\ell^1$.
	
	What follows is the exact same proof as in the previous problem,
	but with $\ell^2$ instead of $\ell^1$.
	Let $\mb x$ be in $\ell^2$ and define $s = \sum_{k=1}^\infty x_k^2$.
	I will show that $\mb x$ is a limit point of $D$
	by proving that it is the limit of the sequence $(\mb x_n)$,
	where $\mb x_n$ is the sequence whose $k$th term is $x_k$
	when $k\le n$, and $0$ when $k > 0$
	(for example $\mb x_3 = (x_1, x_2, x_3, 0, 0, \dots)$).
	We see that
	\[d(\mb x, \mb x_n) = \sum_{k=1}^n (x_k - x_k)^2 +
	\sum_{k=n+1}^\infty (x_k - 0)^2 = \sum_{k=n+1}^\infty x_k^2 =
	s - \sum_{k=1}^n x_k^2 = s - s_n,\]
	where we define $s_n$ to be the partial sums of the sequence $(x_k^2)$.
	Taking the limit, we see that
	\[\lim_{n\to\infty} d(\mb x, \mb x_n) =
	\lim_{n\to\infty} s - s_n = s - \lim_{n\to\infty} s_n = s - s = 0,\]
	so $d(\mb x, \mb x_n)$ will eventually be smaller than $\eps$
	for any choice of $\eps > 0$; thus, $(\mb x_n)$ converges to $\mb x$
	in the metric space $(\ell^2, d)$.
\end{soln}

\begin{problem}
	Let $(X, d)$ be a metric space and $D\subseteq X$ a dense set.
	Suppose that $f\colon X\to\RR$ is zero at all points of $D$.
	Show that if $f$ is continuous
	then $f$ is zero at all points of $X$.
\end{problem}

\begin{soln}
	Since $D$ is a dense subset of $X$,
	every point of $X$ is a limit point of $D$.
	This means that for each $x\in X$, there is a sequence
	$(d_n)$ of elements of $D$ that converges to $x$.
	By continuity, we thus have
	\[f(x) = f\left(\lim_{n\to\infty} d_n\right) =
	\lim_{n\to\infty} f(d_n) = \lim_{n\to\infty} 0 = 0,\]
	so $f(x) = 0$ for all $x\in X$.
\end{soln}

\end{document}